\chapter{Differentiation and Radon--Nikodym}

\begin{center}\itshape
A staircase climbs from $0$ to $1$ without ever jumping, and yet it is flat almost everywhere: its derivative is $0$ at almost
every point. Where did the climbing happen? And when can a function be recovered by integrating its derivative?
\end{center}

In 1930 G.~H.~Hardy and J.~E.~Littlewood published a paper on a problem of complex analysis, and opened it with cricket. A
batsman's scores in a season form a sequence; his ``average'' after each innings depends on where one starts counting. A
batsman who wants to look good will start counting at the innings that makes his average largest. How much can this
flattery inflate the averages? Hardy and Littlewood proved that it is bounded, in a precise sense, and their \emph{maximal
inequality} turned out to be the key to a much older question: Lebesgue's.

Lebesgue had asked in 1904 when the fundamental theorem of calculus holds for his integral: when is $F(b)-F(a)=\int_a^bF'$?
It is not obvious. Vito Volterra had built in 1881 a function that is differentiable everywhere, with a bounded derivative that
is not Riemann integrable. And the devil's staircase of Chapter~2 rises from $0$ to $1$ while its derivative vanishes almost
everywhere: integrating the derivative gives $0$. Lebesgue proved that every monotone function is differentiable almost
everywhere, and that the functions for which the fundamental theorem holds are exactly the \emph{absolutely continuous} ones.

The same story has a version for measures, and it is the one probability needs. When does a measure $\nu$ have a density with
respect to another measure $\mu$, $\nu(A)=\int_Af\,d\mu$? Johann Radon (1913) answered for $\R^d$, Otto Nikodym (1930) in general:
exactly when $\nu$ gives no mass to the sets that $\mu$ ignores. The density $f=\frac{d\nu}{d\mu}$ is a derivative of one measure
against another, and the Radon--Nikodym theorem is the foundation of conditional expectations (Tome~10), of likelihood ratios in
statistics, and of importance sampling in computing. Along the way, every law on $\R$ splits into three parts, atoms, a density
and a singular continuous part, as Tome~3's notebooks announced.

\section{Signed measures and the Hahn decomposition}

A \emph{finite signed measure} on $(E,\mathcal A)$ is a map $\nu:\mathcal A\to\R$ with $\nu(\varnothing)=0$ and $\nu(\bigsqcup_kA_k)=\sum_k\nu(A_k)$ for disjoint $A_k$; the series then converges absolutely, since
its sum does not depend on the order (Tome~0). The difference of two finite measures is one; the Jordan decomposition below says that there are no others. A set $P$ is
\emph{positive} for $\nu$ if $\nu(A)\ge0$ for every measurable $A\subseteq P$, and \emph{negative} if $\nu(A)\le0$ for all such $A$.

\begin{lemma}\label{lem:positivepart}
If $\nu(B)>0$, then $B$ contains a positive set $P$ with $\nu(P)\ge\nu(B)$.
\end{lemma}
\begin{proof}
Remove the worst negative pieces greedily. If $B$ is not positive, let $n_1$ be the smallest integer such that some $C_1\subseteq B$ has $\nu(C_1)<-\frac1{n_1}$. If $B\setminus C_1$ is not positive, let $n_2$ be the
smallest integer such that some $C_2\subseteq B\setminus C_1$ has $\nu(C_2)<-\frac1{n_2}$, and so on (if the process stops, the remaining set is positive and we are done). Let $P=B\setminus\bigcup_kC_k$.
Then $\nu(B)=\nu(P)+\sum_k\nu(C_k)$, and since the $C_k$ are disjoint the series converges, so $\sum_k\frac1{n_k}<\infty$ and $n_k\to\infty$; also $\nu(P)\ge\nu(B)$. If some $A\subseteq P$ had
$\nu(A)<0$, then $\nu(A)<-\frac1{n_k-1}$ for $k$ large, and $A\subseteq B\setminus(C_1\cup\cdots\cup C_{k-1})$ would contradict the minimality of $n_k$. So $P$ is positive.
\end{proof}

\begin{theorem}[Hahn decomposition]\label{thm:hahn}
For a finite signed measure $\nu$ there is a partition $E=P\sqcup N$ with $P$ positive and $N$ negative. If $P'\sqcup N'$ is another, $\nu$ vanishes on all subsets of $P\,\triangle\,P'$.
\end{theorem}
\begin{proof}
Let $m=\sup\{\nu(Q):Q\text{ positive}\}$, finite since $\nu$ is real-valued, and $Q_j$ positive with $\nu(Q_j)\to m$. A countable union of positive sets is positive (cut a subset of the union
into disjoint pieces, each inside one $Q_j$), so $P=\bigcup_jQ_j$ is positive with $\nu(P)=m$. Let $N=E\setminus P$. If some $B\subseteq N$ had $\nu(B)>0$, Lemma~\ref{lem:positivepart} would give a positive
$P''\subseteq B$ with $\nu(P'')>0$, and $P\cup P''$ would be positive with $\nu(P\cup P'')>m$. So $N$ is negative. For uniqueness, a subset of $P\setminus P'=P\cap N'$ is both positive and negative.
\end{proof}

\begin{corollary}[Jordan decomposition]\label{cor:jordan}
Every finite signed measure is $\nu=\nu^+-\nu^-$ with $\nu^\pm$ finite measures, $\nu^+(A)=\nu(A\cap P)$, $\nu^-(A)=-\nu(A\cap N)$, living on the disjoint sets $P$ and $N$. The measure
$|\nu|=\nu^++\nu^-$ is the \emph{total variation} of $\nu$.
\end{corollary}

\section{The Radon--Nikodym theorem}

Two measures $\mu,\nu$ on $(E,\mathcal A)$ are \emph{mutually singular}, $\nu\perp\mu$, if they live on disjoint sets: $E=S\sqcup T$ with $\mu(S)=0$ and $\nu(T)=0$. And $\nu$ is \emph{absolutely continuous} with respect to
$\mu$, $\nu\ll\mu$, if $\mu(A)=0$ implies $\nu(A)=0$. If $\nu(A)=\int_Af\,d\mu$ for a measurable $f\ge0$, then clearly $\nu\ll\mu$. The theorem says that the converse holds.

\begin{think}
Let $\mu$ be the law of a standard Gaussian and $\nu$ the law of a Gaussian of mean $1$ and variance $1$. Is there a function $f$ with $\nu(A)=\int_Af\,d\mu$ for every Borel $A$? What about $\nu$ the law
of a Poisson variable?
\end{think}

For Gaussians, $f(x)=\frac{e^{-(x-1)^2/2}}{e^{-x^2/2}}=e^{x-\frac12}$, the ratio of the densities. For a Poisson law $\nu$, which lives on the null set $\N$, no such $f$ exists: $\nu(\N)=1$ but $\mu(\N)=0$. The
Poisson law and the Gaussian are mutually singular.

\begin{theorem}[Radon--Nikodym]\label{thm:rn}
Let $\mu$ and $\nu$ be $\sigma$-finite measures with $\nu\ll\mu$. There is a measurable $f\ge0$ with $\nu(A)=\int_Af\,d\mu$ for all $A\in\mathcal A$, unique $\mu$-a.e. It is written $\frac{d\nu}{d\mu}$, the
\emph{Radon--Nikodym derivative}.
\end{theorem}
\begin{proof}
\emph{Finite measures.} Let $\mathcal G$ be the set of measurable $g\ge0$ with $\int_Ag\,d\mu\le\nu(A)$ for all $A$, and $m=\sup_{g\in\mathcal G}\int g\,d\mu\le\nu(E)$. If $g,h\in\mathcal G$, then $\max(g,h)\in\mathcal G$ (split $A$
into $A\cap\{g\ge h\}$ and $A\cap\{g<h\}$). So there are $g_n\in\mathcal G$, which we may take increasing, with $\int g_n\to m$; by monotone convergence $f=\lim g_n\in\mathcal G$ and $\int f=m$. Then
$\nu_0(A)=\nu(A)-\int_Af\,d\mu$ is a finite measure, and we claim $\nu_0=0$. If not, $\nu_0(E)>0$, so $\mu(E)>0$ (as $\nu_0\le\nu\ll\mu$), and we can choose $\varepsilon>0$ with $\nu_0(E)>\varepsilon\mu(E)$. Let $P\sqcup N$ be a
Hahn decomposition for $\nu_0-\varepsilon\mu$. Then $(\nu_0-\varepsilon\mu)(P)\ge(\nu_0-\varepsilon\mu)(E)>0$, so $\mu(P)>0$ (otherwise $\nu_0(P)=0$ too). For every $A$,
\[\int_A(f+\varepsilon\mathbf 1_P)\,d\mu=\int_Af\,d\mu+\varepsilon\mu(A\cap P)\le\int_Af\,d\mu+\nu_0(A\cap P)\le\nu(A),\]
since $P$ is positive for $\nu_0-\varepsilon\mu$. So $f+\varepsilon\mathbf 1_P\in\mathcal G$, with integral $m+\varepsilon\mu(P)>m$: a contradiction.

\emph{$\sigma$-finite measures.} Cut $E$ into disjoint $E_n$ with $\mu(E_n)$ and $\nu(E_n)$ finite (intersect the pieces for $\mu$ and for $\nu$), apply the finite case on each $E_n$, and glue the densities.

\emph{Uniqueness.} If $\int_Af=\int_Ag$ for all $A$, take $A=\{f>g\}\cap E_n\cap\{f\le k\}$: then $\int_A(f-g)=0$ with $f-g>0$ on $A$, so $\mu(A)=0$; let $k,n\to\infty$, and exchange $f,g$.
\end{proof}

\begin{theorem}[Lebesgue decomposition]\label{thm:lebdecomp}
Let $\mu,\nu$ be $\sigma$-finite. Then $\nu=\nu_a+\nu_s$ uniquely with $\nu_a\ll\mu$ and $\nu_s\perp\mu$. Hence $\nu=f\mu+\nu_s$ with $f=\frac{d\nu_a}{d\mu}$.
\end{theorem}
\begin{proof}
Assume $\nu$ finite (the $\sigma$-finite case follows piece by piece). Let $s=\sup\{\nu(Z):\mu(Z)=0\}$, and $Z_j$ $\mu$-null with $\nu(Z_j)\to s$. Then $Z=\bigcup_jZ_j$ is $\mu$-null and $\nu(Z)=s$. Put
$\nu_s=\nu(\cdot\cap Z)$, which lives on the $\mu$-null set $Z$, and $\nu_a=\nu(\cdot\setminus Z)$. If $\mu(A)=0$, then $Z\cup A$ is $\mu$-null, so $s\ge\nu(Z\cup A)=s+\nu_a(A)$ and $\nu_a(A)=0$: $\nu_a\ll\mu$. Uniqueness: if
$\nu_a+\nu_s=\nu_a'+\nu_s'$, the finite signed measure $\nu_a-\nu_a'=\nu_s'-\nu_s$ vanishes on $\mu$-null sets and lives on a $\mu$-null set (the union of those of $\nu_s,\nu_s'$), so it is $0$.
\end{proof}

\begin{corollary}[the three parts of a law]\label{cor:threeparts}
Every probability $\mu$ on $\R$ is uniquely $\mu=\mu_d+\mu_{ac}+\mu_{sc}$, where $\mu_d=\sum_ip_i\delta_{a_i}$ is carried by countably many atoms, $\mu_{ac}$ has a density with respect to $\lambda$, and $\mu_{sc}$ has no atoms and
lives on a $\lambda$-null set.
\end{corollary}
\begin{proof}
The atoms $\{a:\mu(\{a\})>0\}$ are countable: at most $n$ of them have mass $>\frac1n$. Let $\mu_d$ be $\mu$ restricted to them, and decompose the rest, which has no atoms, with Theorem~\ref{thm:lebdecomp}
with respect to $\lambda$; restrictions of an atomless measure are atomless.
\end{proof}

\begin{figure}[ht]
\centering
\fig{ch10-threeparts}
\caption{The distribution function of a law with three parts: weight $\frac14$ on an atom at $-1$ (a jump), weight $\frac12$ with a density on $[0,1]$ (a smooth rise), and weight $\frac14$ spread on a Cantor set
in $[2,3]$ (a devil's staircase, continuous, but flat almost everywhere).}
\label{fig:threeparts}
\end{figure}

Exercise~\ref{ex:salem} shows that a singular continuous law can even have a strictly increasing distribution function. In probability, if $\mathbb Q\ll\mathbb P$, the density $L=\frac{d\mathbb Q}{d\mathbb P}$ is
the \emph{likelihood ratio}: $\mathbb E_{\mathbb Q}[X]=\mathbb E_{\mathbb P}[XL]$. This is the principle of importance sampling (this chapter's notebook), of the Neyman--Pearson lemma in statistics, and of
Girsanov's theorem in Tome~11. Exercise~\ref{ex:condexp} builds conditional expectations from Theorem~\ref{thm:rn}.

\section{The maximal function and Lebesgue's differentiation theorem}

On $\R^d$, a density can be recovered from its measure by zooming in: if $\nu=f\lambda_d$ and $f$ is continuous at $x$, then $\frac{\nu(B(x,r))}{\lambda_d(B(x,r))}\to f(x)$ as $r\to0$. Lebesgue's theorem says that this
works almost everywhere for every integrable $f$. The difficulty is to control all radii at once, and that is the job of the maximal function. A measurable $f$ is \emph{locally integrable}, $f\in L^1_{\mathrm{loc}}$, if
$\int_K|f|<\infty$ for every compact $K$. Balls are open balls for the Euclidean norm, and $B(x,r)$ has measure $\omega_dr^d$.

\begin{lemma}[Vitali's covering lemma]\label{lem:vitalicover}
From finitely many open balls $B_1,\dots,B_n$ one can extract disjoint ones $B_{i_1},\dots,B_{i_k}$ with $\lambda_d(B_1\cup\cdots\cup B_n)\le3^d\sum_j\lambda_d(B_{i_j})$.
\end{lemma}
\begin{proof}
Greedily: take the largest ball, discard every ball meeting it, take the largest remaining ball, and so on. A discarded ball meets a chosen ball at least as large, so it lies inside the concentric
ball with three times the radius, of measure $3^d$ times larger. The tripled chosen balls cover everything.
\end{proof}

For $f\in L^1_{\mathrm{loc}}(\R^d)$, the \emph{Hardy--Littlewood maximal function} is
\[Mf(x)=\sup_{r>0}\frac1{\lambda_d(B(x,r))}\int_{B(x,r)}|f|\,d\lambda_d.\]
It is the batsman's flattering average. $\{Mf>t\}$ is open: if $\frac1{\lambda_d(B(x,r))}\int_{B(x,r)}|f|>t$, the same holds for $B(x',r+|x-x'|)$ when $x'$ is close to $x$, by continuity in the radius. So $Mf$ is
Borel.

\begin{theorem}[maximal inequality]\label{thm:maximal}
For $f\in L^1(\R^d)$ and $t>0$, $\lambda_d(\{Mf>t\})\le\dfrac{3^d}t\|f\|_1$.
\end{theorem}
\begin{proof}
Let $K\subseteq\{Mf>t\}$ be compact. Each $x\in K$ is the centre of a ball $B_x$ with $\int_{B_x}|f|>t\lambda_d(B_x)$; finitely many cover $K$, and Lemma~\ref{lem:vitalicover} extracts disjoint ones with
\[\lambda_d(K)\le3^d\sum_j\lambda_d(B_{i_j})<\frac{3^d}t\sum_j\int_{B_{i_j}}|f|\le\frac{3^d}t\|f\|_1.\]
Take the supremum over $K$ (inner regularity, Proposition~\ref{prop:regularity}).
\end{proof}

\begin{figure}[ht]
\centering
\fig{ch10-maximal}
\caption{The indicator of $[-1,1]$ (blue) and its maximal function (orange), equal to $1$ on $[-1,1]$ and to $\frac1{|x|+1}$ outside. The maximal function of an integrable function is usually not integrable,
because of the tail $\frac1{|x|}$, but the measure of $\{Mf>t\}$ is of order $\frac1t$: the ``weak'' bound of Theorem~\ref{thm:maximal}.}
\label{fig:maximal}
\end{figure}

\begin{theorem}[Lebesgue's differentiation theorem]\label{thm:lebdiff}
Let $f\in L^1_{\mathrm{loc}}(\R^d)$. For almost every $x$,
\[\lim_{r\to0}\frac1{\lambda_d(B(x,r))}\int_{B(x,r)}|f(y)-f(x)|\,dy=0.\]
Such $x$ are the \emph{Lebesgue points} of $f$; at each of them the averages of $f$ over $B(x,r)$ tend to $f(x)$.
\end{theorem}
\begin{proof}
The statement is local, so we may assume $f\in L^1$ (multiply by the indicator of a large ball). Let $\Omega f(x)=\limsup_{r\to0}\frac1{\lambda_d(B(x,r))}\int_{B(x,r)}|f(y)-f(x)|\,dy$. If $g$ is continuous, $\Omega g=0$
everywhere. For $g\in C_c$, since $|f(y)-f(x)|\le|(f-g)(y)|+|g(y)-g(x)|+|(g-f)(x)|$, we get $\Omega f\le M(f-g)+|f-g|$. Hence, for $t>0$,
\[\lambda_d(\{\Omega f>2t\})\le\lambda_d(\{M(f-g)>t\})+\lambda_d(\{|f-g|>t\})\le\frac{3^d+1}t\|f-g\|_1\]
(Theorem~\ref{thm:maximal} and Markov). By Theorem~\ref{thm:ccdense} the right side can be made as small as we wish, so $\lambda_d^*(\{\Omega f>2t\})=0$ for every $t>0$, and $\Omega f=0$ a.e.
\end{proof}

In dimension one, intervals $[x,x+h]$ and $[x-h,x]$ lie in $B(x,|h|)$ and have half its length, so at every Lebesgue point $\frac1h\int_x^{x+h}f\to f(x)$ as $h\to0$ from either side: $F(x)=\int_a^xf$ is
differentiable at $x$ with $F'(x)=f(x)$. This is half of the fundamental theorem of calculus, now for every integrable $f$, at almost every point.

\section{Differentiating measures and monotone functions}

A Borel measure on $\R^d$ that is finite on bounded sets is \emph{locally finite}. To differentiate such measures, we need one more regularity fact; Chapter~11 proves it in every metric space.

\begin{lemma}\label{lem:outerreg}
A finite Borel measure $\nu$ on a metric space $X$ is outer regular: $\nu(A)=\inf\{\nu(U):U\supseteq A\text{ open}\}$ for every Borel set $A$.
\end{lemma}
\begin{proof}
Let $\mathcal C$ be the family of Borel sets $A$ such that for every $\varepsilon>0$ there are a closed $C$ and an open $U$ with $C\subseteq A\subseteq U$ and $\nu(U\setminus C)<\varepsilon$. A closed set $C$ is in $\mathcal C$: the open
sets $U_n=\{x:d(x,C)<\frac1n\}$ decrease to $C$, so $\nu(U_n\setminus C)\to0$ by continuity from above. $\mathcal C$ is stable under complements (exchange the roles of $C$ and $U$) and countable unions: for $A=\bigcup_kA_k$,
take $C_k\subseteq A_k\subseteq U_k$ with $\nu(U_k\setminus C_k)<\varepsilon2^{-k}$, let $U=\bigcup_kU_k$ and $C=C_1\cup\cdots\cup C_K$ with $K$ so large that $\nu(\bigcup_kC_k\setminus C)<\varepsilon$. So $\mathcal C$ is a $\sigma$-algebra
containing the closed sets, hence all Borel sets.
\end{proof}

\begin{theorem}[differentiation of measures]\label{thm:diffmeasures}
Let $\nu$ be a locally finite Borel measure on $\R^d$, with Lebesgue decomposition $\nu=f\lambda_d+\nu_s$. Then for almost every $x$,
\[\lim_{r\to0}\frac{\nu(B(x,r))}{\lambda_d(B(x,r))}=f(x).\]
\end{theorem}
\begin{proof}
By Theorem~\ref{thm:lebdiff} the absolutely continuous part contributes $f(x)$ at almost every point ($f$ is locally integrable as $\nu$ is locally finite). It remains to show that
$\frac{\nu_s(B(x,r))}{\lambda_d(B(x,r))}\to0$ a.e. Working in a large ball, we may assume $\nu_s$ finite. It lives on a $\lambda_d$-null Borel set $Z$: $\nu_s(\R^d\setminus Z)=0$. Fix $k\ge1$ and let
\[F_k=\Bigl\{x\notin Z:\limsup_{r\to0}\frac{\nu_s(B(x,r))}{\lambda_d(B(x,r))}>\frac1k\Bigr\}.\]
Given $\varepsilon>0$, Lemma~\ref{lem:outerreg} gives an open $U\supseteq\R^d\setminus Z$ with $\nu_s(U)<\varepsilon$. Each $x\in F_k$ is the centre of a ball $B_x\subseteq U$ with $\nu_s(B_x)>\frac1k\lambda_d(B_x)$. For a compact
$K\subseteq\bigcup_xB_x$, finitely many $B_x$ cover $K$, and Lemma~\ref{lem:vitalicover} gives disjoint ones with $\lambda_d(K)\le3^d\sum_j\lambda_d(B_{i_j})<3^dk\sum_j\nu_s(B_{i_j})\le3^dk\varepsilon$. By inner regularity
$\lambda_d(\bigcup_xB_x)\le3^dk\varepsilon$, so $\lambda_d^*(F_k)\le3^dk\varepsilon$ for every $\varepsilon$: $F_k$ is null. Since $Z$ is null too, the $\limsup$ vanishes almost everywhere.
\end{proof}

\begin{think}
Can a continuous increasing function have derivative $0$ almost everywhere and still rise from $0$ to $1$?
\end{think}

Yes: the devil's staircase $c$ of Exercise~\ref{ex:devil} is constant on each interval removed to build the Cantor set, so $c'=0$ on a set of measure $1$, and yet $c(1)-c(0)=1$. Its measure $\mu_c$ lives on the
Cantor set: it is singular, and Theorem~\ref{thm:diffmeasures} sees only its density, $0$. All the climbing happens on a set of measure zero.

\begin{theorem}[Lebesgue, 1904]\label{thm:monotone}
A monotone function $F:\R\to\R$ is differentiable almost everywhere. If $F$ is non-decreasing and right-continuous, then $F'=\frac{d\mu_{F,a}}{d\lambda}$ a.e., where $\mu_{F,a}$ is the absolutely continuous part of
$\mu_F$ (Theorem~\ref{thm:stieltjes}), and $\int_a^bF'\le F(b)-F(a)$.
\end{theorem}
\begin{proof}
Let $F$ be non-decreasing and right-continuous, and $\mu_F=f\lambda+\nu_s$. For $h>0$, $F(x+h)-F(x)=\mu_F(]x,x+h])$ and $F(x)-F(x-h)=\mu_F(]x-h,x])$, intervals of length $h$ inside $B(x,2h)$, which has length $4h$. At every $x$ which is a Lebesgue point of $f$ and where $\frac{\nu_s(B(x,r))}{2r}\to0$ (almost every $x$, by Theorems~\ref{thm:lebdiff} and~\ref{thm:diffmeasures}),
\[\Bigl|\frac{F(x+h)-F(x)}h-f(x)\Bigr|\le\frac1h\int_x^{x+h}|f-f(x)|+\frac{\nu_s(B(x,2h))}h\to0,\]
and similarly from the left; so $F'(x)=f(x)$. Then $\int_a^bF'=\int_{]a,b]}f\,d\lambda\le\mu_F(]a,b])=F(b)-F(a)$.

For a general non-decreasing $F$, let $G(x)=F(x^+)$, which is non-decreasing and right-continuous. The function $H=G-F\ge0$ vanishes except at the discontinuities $x_j$ of $F$, a countable set, and
$\sum_{x_j\in[-n,n]}H(x_j)\le F(n+1)-F(-n-1)$ for each $n$ (the jumps are disjoint parts of the rise). So $\rho=\sum_jH(x_j)\delta_{x_j}$ is a locally finite measure, singular with respect to $\lambda$. At every $x$
outside $\{x_j\}$ with $\frac{\rho(B(x,r))}{2r}\to0$ (almost every $x$, by Theorem~\ref{thm:diffmeasures}), $H(x)=0$ and $0\le\frac{H(x+h)}{|h|}\le\frac{\rho(B(x,2|h|))}{|h|}\to0$: so $H'(x)=0$, and $F'=G'$ a.e. A non-increasing $F$ is
$-(-F)$.
\end{proof}

\section{Absolutely continuous functions}

\begin{definition}
$F:[a,b]\to\R$ is \emph{absolutely continuous} if for every $\varepsilon>0$ there is $\delta>0$ such that $\sum_i|F(b_i)-F(a_i)|<\varepsilon$ for every finite family of disjoint intervals $]a_i,b_i[$ of total length
$\sum_i(b_i-a_i)<\delta$.
\end{definition}

With one interval, this is uniform continuity; with many, it forbids the devil's staircase, which rises by $1$ on the Cantor set, covered by intervals of arbitrarily small total length.

\begin{theorem}[fundamental theorem of calculus]\label{thm:ftc}
For $F:[a,b]\to\R$, the following are equivalent: (i) $F$ is absolutely continuous; (ii) there is $f\in L^1([a,b])$ with $F(x)=F(a)+\int_a^xf$ for all $x$. Then $F$ is differentiable a.e., $F'=f$ a.e., and
$F(b)-F(a)=\int_a^bF'$.
\end{theorem}
\begin{proof}
(ii)$\Rightarrow$(i): by Lemma~\ref{lem:abscont}, $\int_A|f|<\varepsilon$ when $\lambda(A)<\delta$; take $A=\bigcup_i\,]a_i,b_i[$. And $F'=f$ a.e. by Lebesgue's differentiation theorem.

(i)$\Rightarrow$(ii). \emph{Bounded variation.} Let $\delta$ correspond to $\varepsilon=1$. For a subdivision $P:a=x_0<\cdots<x_n=x$, let $V(F,P)=\sum_i|F(x_i)-F(x_{i-1})|$ and $T(x)=\sup_PV(F,P)$, the variation of $F$ on $[a,x]$.
Cutting $[a,b]$ into $m$ pieces of length less than $\delta$, and refining any $P$ by their endpoints, shows $T(b)\le m<\infty$. For $x<y$, $T(y)\ge T(x)+|F(y)-F(x)|$ (append $y$ to the subdivisions of $[a,x]$), so
$T$ and $T-F$ are non-decreasing.

\emph{$T$ is absolutely continuous.} Given $\varepsilon$ and its $\delta$, let $]a_i,b_i[$ be disjoint with total length $<\delta$. For subdivisions $P_i$ of the $[a_i,b_i]$, all their small intervals together have total length
$<\delta$, so $\sum_iV(F|_{[a_i,b_i]},P_i)<\varepsilon$. Taking suprema, $\sum_i(T(b_i)-T(a_i))\le\varepsilon$, since $T(b_i)-T(a_i)$ is the variation of $F$ on $[a_i,b_i]$ (subdivisions of $[a,b_i]$ may be assumed to contain $a_i$).
So $T$ and $T-F$ are non-decreasing and absolutely continuous, in particular continuous.

\emph{Non-decreasing absolutely continuous $G$.} Extend $G$ by constants outside $[a,b]$; then $\mu_G\ll\lambda$. Indeed, let $\lambda(N)=0$ and $\varepsilon>0$; cover $N\cap\,]a,b[$ by open intervals of total length $<\delta$, which we may
take disjoint (a union of open intervals is a disjoint union of open intervals, of no larger total length) and contained in $]a,b[$; every finite subfamily has $\sum(G(b_i)-G(a_i))<\varepsilon$, so $\mu_G(N)\le\varepsilon$ (the points
$a,b$ carry no mass, $G$ being continuous). By Radon--Nikodym, $\mu_G=g\lambda$, and $G(x)-G(a)=\mu_G(]a,x])=\int_a^xg$. Applying this to $T$ and to $T-F$ gives (ii), with $f$ the density of $T$ minus the density of $T-F$.
\end{proof}

The devil's staircase is continuous, non-decreasing, differentiable almost everywhere with integrable derivative $0$, and not absolutely continuous: absolute continuity is exactly what is needed
for the fundamental theorem of calculus, no more and no less.

\begin{beauty}
Three ideas that look unrelated turn out to be one. The density of a law, the derivative of a function, and the limit of averages over shrinking balls are the same object seen from three sides:
$\frac{d\nu}{d\lambda}$. And a story about a batsman flattering his cricket average controls them all.
\end{beauty}

\section{Exercises}

\begin{exercise}[$\star$]\label{ex:cantorfn}
Let $c$ be the devil's staircase (Exercise~\ref{ex:devil}). Show that $c'=0$ on $[0,1]\setminus C$, that $c$ is not absolutely continuous, and that its Lebesgue--Stieltjes measure $\mu_c$ is the Cantor law, singular
and without atoms. What is the distribution function of $\mu_c$ restricted to $[0,\frac13]$, in terms of $c$?
\end{exercise}

\begin{exercise}[$\star\star$]\label{ex:condexp}
\emph{Conditional expectation (Tome~3, Chapter~10).} Let $X\in L^1(\Omega,\mathcal F,\mathbb P)$ and $\mathcal G\subseteq\mathcal F$ a $\sigma$-algebra. (a) For $X\ge0$, show that $\nu(A)=\mathbb E[X\mathbf 1_A]$, $A\in\mathcal G$, is a finite
measure on $\mathcal G$ with $\nu\ll\mathbb P|_{\mathcal G}$, and let $\mathbb E[X\mid\mathcal G]=\frac{d\nu}{d\mathbb P|_{\mathcal G}}$. (b) Extend to integrable $X$ and show that $Y=\mathbb E[X\mid\mathcal G]$ is the unique $\mathcal G$-measurable
integrable $Y$ with $\mathbb E[Y\mathbf 1_A]=\mathbb E[X\mathbf 1_A]$ for all $A\in\mathcal G$. (c) Prove linearity, $\mathbb E[\mathbb E[X\mid\mathcal G]]=\mathbb E[X]$, and $\mathbb E[ZX\mid\mathcal G]=Z\,\mathbb E[X\mid\mathcal G]$ for bounded
$\mathcal G$-measurable $Z$.
\end{exercise}

\begin{exercise}[$\star$]\label{ex:likelihood}
Let $\mathbb P=\mathcal N(0,1)$ and $\mathbb Q=\mathcal N(m,1)$ on $\R$. Compute $L=\frac{d\mathbb Q}{d\mathbb P}$, check $\mathbb E_{\mathbb P}[L]=1$, and compute the relative entropy $\int\ln L\,d\mathbb Q$. Show that $\mathcal N(0,1)$ and
$\mathcal N(0,2)$ are mutually absolutely continuous, and that $\mathcal N(0,1)$ and a Poisson law are mutually singular.
\end{exercise}

\begin{exercise}[$\star\star\star$]\label{ex:salem}
\emph{A strictly increasing singular function.} Let $0<p<1$, $p\ne\frac12$, let $B_1,B_2,\ldots$ be independent with $\mathbb P(B_k=1)=p=1-\mathbb P(B_k=0)$, and $F_p$ the distribution function of $X=\sum_kB_k2^{-k}$.
Show that $F_p$ is continuous and strictly increasing on $[0,1]$, and that $F_p'=0$ almost everywhere. (Use the strong law of large numbers, Tome~3, and Borel's theorem of Chapter~6.)
\end{exercise}

\begin{exercise}[$\star\star$]\label{ex:density}
\emph{Lebesgue's density theorem.} For a Borel set $A\subseteq\R^d$, show that $\frac{\lambda_d(A\cap B(x,r))}{\lambda_d(B(x,r))}\to\mathbf 1_A(x)$ for almost every $x$. Deduce that there is no Borel $A\subseteq[0,1]$ with
$\lambda(A\cap I)=\frac12\lambda(I)$ for every interval $I\subseteq[0,1]$. Compare with Exercise~\ref{ex:everywhere}.
\end{exercise}

\begin{exercise}[$\star\star$]\label{ex:variation}
(a) Show that $x\sin\frac1x$ (with value $0$ at $0$) is continuous on $[0,1]$ but not of bounded variation. (b) Show that $F(x)=x^2\sin\frac1{x^2}$ (with $F(0)=0$) is differentiable at every point of $[-1,1]$, but that
$F'$ is not integrable on $[0,1]$. Is $F$ absolutely continuous on $[0,1]$?
\end{exercise}

\section{Hints}

\begin{hint}{ex:cantorfn}
A continuous function constant on an interval has derivative $0$ there. For absolute continuity, cover $C_n$ by $2^n$ intervals of total length $(\frac23)^n$. On $[0,\frac13]$, $c(x)=\frac12c(3x)$.
\end{hint}
\begin{hint}{ex:condexp}
(b) Split $X=X^+-X^-$. Uniqueness: apply the defining property to $A=\{Y>Y'\}$. (c) For $Z=\mathbf 1_B$, $B\in\mathcal G$, the property is the definition; then simple $Z$, then limits.
\end{hint}
\begin{hint}{ex:likelihood}
$L$ is the ratio of densities. $\ln L(x)=mx-\frac{m^2}2$.
\end{hint}
\begin{hint}{ex:salem}
Each dyadic interval of rank $n$ has probability $p^k(1-p)^{n-k}>0$. The set of $x$ whose binary digits have frequency of ones equal to $p$ has $\mu_p$-measure $1$ and Lebesgue measure $0$. Then apply
Theorem~\ref{thm:monotone}.
\end{hint}
\begin{hint}{ex:density}
Apply Theorem~\ref{thm:lebdiff} to $\mathbf 1_{A\cap B(0,R)}$. For the second part, the density at almost every point of $[0,1]$ would be $\frac12$.
\end{hint}
\begin{hint}{ex:variation}
(a) Look at the points $x_k=\frac1{(k+\frac12)\pi}$. (b) $F'(x)=2x\sin\frac1{x^2}-\frac2x\cos\frac1{x^2}$ for $x\ne0$; integrate $|F'|$ over the intervals where $|\cos\frac1{x^2}|\ge\frac12$.
\end{hint}

\section{Solutions}

\begin{solution}{ex:cantorfn}
\idea Everything is visible on the staircase: it is flat on the gaps of the Cantor set, which have total length $1$, and it climbs only on the Cantor set, which can be covered by intervals of tiny total
length; that is exactly what absolute continuity forbids.
\writeup By Exercise~\ref{ex:devil}, $c$ is constant on each removed open interval, so $c'=0$ on $[0,1]\setminus C$, a set of measure $1$. The stage $C_n$ consists of $2^n$ closed intervals of length $3^{-n}$; on each,
$c$ increases by $2^{-n}$ (by self-similarity, $c(x)=\frac12c(3x)$ on $[0,\frac13]$ and $c(x)=\frac12+\frac12c(3x-2)$ on $[\frac23,1]$). Their total length $(\frac23)^n$ tends to $0$, while $\sum|c(b_i)-c(a_i)|=1$: $c$ is not absolutely
continuous (slightly enlarge the intervals to open ones of total length less than $\delta$). Extend $c$ by $0$ on $]-\infty,0[$ and $1$ on $]1,\infty[$. Then $\mu_c$ is a probability with $\mu_c([0,1]\setminus C)=0$ (it is $0$
on each gap, where $c$ is constant), $\lambda(C)=0$, and $\mu_c(\{x\})=0$ by continuity: it is singular and atomless, and it is the law of $\sum_k2B_k3^{-k}$ with fair bits $B_k$ (both have distribution function $c$,
by self-similarity). Restricted to $[0,\frac13]$, $\mu_c$ has distribution function $c(x)=\frac12c(3x)$: it is half the image of $\mu_c$ under $x\mapsto\frac x3$.
\end{solution}

\begin{solution}{ex:condexp}
\idea Conditional expectation given $\mathcal G$ is a density: the measure ``$X$ restricted to the information $\mathcal G$'' has a density with respect to $\mathbb P$ on $\mathcal G$, and the density must be
$\mathcal G$-measurable because it is a function on the smaller $\sigma$-algebra. Radon--Nikodym does all the work.
\writeup (a) $\nu$ is countably additive by monotone convergence, finite as $\mathbb E[X]<\infty$, and $\mathbb P(A)=0$ implies $\mathbb E[X\mathbf 1_A]=0$. Theorem~\ref{thm:rn} on $(\Omega,\mathcal G)$ gives a $\mathcal G$-measurable $Y\ge0$ with
$\mathbb E[Y\mathbf 1_A]=\nu(A)$, unique a.s. (b) Set $\mathbb E[X\mid\mathcal G]=\mathbb E[X^+\mid\mathcal G]-\mathbb E[X^-\mid\mathcal G]$, integrable since $\mathbb E[\mathbb E[X^\pm\mid\mathcal G]]=\mathbb E[X^\pm]$. If $Y,Y'$ both satisfy the property,
$A=\{Y>Y'\}\in\mathcal G$ gives $\mathbb E[(Y-Y')\mathbf 1_A]=0$ with $Y-Y'>0$ on $A$, so $\mathbb P(A)=0$; similarly $\{Y'>Y\}$. (c) Linearity follows from uniqueness: $a\mathbb E[X\mid\mathcal G]+b\mathbb E[X'\mid\mathcal G]$ has the defining
property of $\mathbb E[aX+bX'\mid\mathcal G]$. Taking $A=\Omega$ gives $\mathbb E[\mathbb E[X\mid\mathcal G]]=\mathbb E[X]$. For $Z=\mathbf 1_B$ with $B\in\mathcal G$ and $A\in\mathcal G$, $\mathbb E[\mathbf 1_BY\mathbf 1_A]=\mathbb E[Y\mathbf 1_{A\cap B}]=\mathbb E[X\mathbf 1_{A\cap B}]=\mathbb E[\mathbf 1_BX\mathbf 1_A]$,
so $\mathbf 1_BY=\mathbb E[\mathbf 1_BX\mid\mathcal G]$. By linearity this holds for simple $Z$; a bounded $\mathcal G$-measurable $Z$ is a uniform limit of simple $Z_n$ bounded by the same constant, and dominated convergence (with
bounds $C|X|$ and $C|Y|$) passes to the limit in $\mathbb E[Z_nY\mathbf 1_A]=\mathbb E[Z_nX\mathbf 1_A]$.
\end{solution}

\begin{solution}{ex:likelihood}
\idea Both laws have densities with respect to Lebesgue measure, positive everywhere, so each is absolutely continuous with respect to the other and the likelihood ratio is the quotient of
densities.
\writeup With $\varphi(x)=\frac{e^{-x^2/2}}{\sqrt{2\pi}}$, $\mathbb Q(A)=\int_A\varphi(x-m)\,dx=\int_A\frac{\varphi(x-m)}{\varphi(x)}\,d\mathbb P(x)$, so $L(x)=e^{mx-m^2/2}$. Then $\mathbb E_{\mathbb P}[L]=\mathbb Q(\R)=1$ (or: $\mathbb E[e^{mZ}]=e^{m^2/2}$), and
$\int\ln L\,d\mathbb Q=\mathbb E_{\mathbb Q}[mX]-\frac{m^2}2=m^2-\frac{m^2}2=\frac{m^2}2$. $\mathcal N(0,1)$ and $\mathcal N(0,2)$ both have positive densities, so the same quotient argument works both ways. A Poisson law lives on the null set
$\N$, which has $\mathcal N(0,1)$-measure $0$: they are mutually singular.
\end{solution}

\begin{solution}{ex:salem}
\idea Under $\mu_p$ the binary digits of $X$ are biased coins, so by the law of large numbers the proportion of ones is $p$; under Lebesgue measure it is $\frac12$. The two measures live on disjoint sets, so
$\mu_p$ is singular, yet it charges every dyadic interval, so its distribution function climbs everywhere.
\writeup Let $\mu_p$ be the law of $X$. \emph{No atoms.} A point has at most two binary expansions, and each expansion has probability $\lim_np^{k_n}(1-p)^{n-k_n}\le\lim_n\max(p,1-p)^n=0$. So $F_p$ is continuous.
\emph{Strictly increasing.} The dyadic interval $[\frac j{2^n},\frac{j+1}{2^n}]$ contains $X$ as soon as the first $n$ bits are those of $j$, an event of probability $p^k(1-p)^{n-k}>0$; every interval $[s,t]\subseteq[0,1]$ with $s<t$ contains
such a dyadic interval, so $F_p(t)>F_p(s)$. \emph{Singular.} Let $S_q$ be the set of $x\in[0,1[$ whose binary digits $D_j(x)$ satisfy $\frac1n\sum_{j\le n}D_j(x)\to q$. By the strong law of large numbers, $\mu_p(S_p)=1$ (the digits of $X$ are the
$B_k$, except on the null event that $X$ is dyadic). By Borel's theorem (Theorem~\ref{thm:borelnormal}), $\lambda(S_{1/2})=1$, and $S_p\cap S_{1/2}=\varnothing$ as $p\ne\frac12$, so $\lambda(S_p)=0$. Thus $\mu_p\perp\lambda$, its absolutely
continuous part is $0$, and Theorem~\ref{thm:monotone} gives $F_p'=0$ almost everywhere. A continuous, strictly increasing function whose derivative is $0$ almost everywhere: the climbing happens on a set of
measure zero which is, this time, dense.
\end{solution}

\begin{solution}{ex:density}
\idea The density of $A$ at $x$ is the average of $\mathbf 1_A$ over a small ball; Lebesgue's theorem says that averages recover the function at almost every point, and $\mathbf 1_A$ takes only the values $0$
and $1$.
\writeup Fix $R$ and $f=\mathbf 1_{A\cap B(0,R)}\in L^1$. By Theorem~\ref{thm:lebdiff}, for a.e. $x$ with $|x|<R$, $\frac{\lambda_d(A\cap B(x,r))}{\lambda_d(B(x,r))}=\frac1{\lambda_d(B(x,r))}\int_{B(x,r)}f\to f(x)=\mathbf 1_A(x)$ for small $r$; let $R$ run
through the integers. If $\lambda(A\cap I)=\frac12\lambda(I)$ for every interval, then the density of $A$ at every point of $]0,1[$ is $\frac12$, which is neither $0$ nor $1$, so the set where the conclusion fails is
all of $]0,1[$, not null: contradiction. Exercise~\ref{ex:everywhere} built a set meeting every interval in positive measure, together with its complement; the density theorem says that the proportions cannot be
the same everywhere.
\end{solution}

\begin{solution}{ex:variation}
\idea A function of bounded variation cannot oscillate with amplitude $\frac1k$ infinitely often near a point, because $\sum\frac1k=\infty$. An everywhere-differentiable function can have a derivative so large
near $0$ that it is not integrable, and then the fundamental theorem fails in Lebesgue's form, so the function is not absolutely continuous.
\writeup (a) Continuity at $0$: $|x\sin\frac1x|\le|x|$. At $x_k=\frac1{(k+\frac12)\pi}$, $x_k\sin\frac1{x_k}=\frac{(-1)^k}{(k+\frac12)\pi}$, so consecutive values differ by more than $\frac1{(k+1)\pi}$, and the variation over
$[0,1]$ is at least $\sum_k\frac1{(k+1)\pi}=\infty$. (b) At $x\ne0$ the formula of the hint holds; at $0$, $|\frac{F(h)}h|\le|h|\to0$, so $F'(0)=0$. On $I_k=[(2k\pi+\frac\pi3)^{-1/2},(2k\pi-\frac\pi3)^{-1/2}]$, where
$\cos\frac1{x^2}\ge\frac12$, we have $|F'(x)|\ge\frac1x-2x\ge\frac1{2x}$ for $x\le\frac12$, so
\[\int_{I_k}|F'|\ge\frac12\ln\frac{(2k\pi+\frac\pi3)^{1/2}}{(2k\pi-\frac\pi3)^{1/2}}=\frac14\ln\frac{6k+1}{6k-1}\ge\frac c k\]
for some $c>0$, and $\sum_k\frac ck=\infty$. So $F'$ is not integrable on $[0,1]$, and by Theorem~\ref{thm:ftc} $F$ is not absolutely continuous there (otherwise $F'$ would be integrable). It is absolutely continuous on every
$[\varepsilon,1]$, where $F'$ is bounded: everything goes wrong at a single point.
\end{solution}

\section*{Sources and further reading}

\begin{itemize}
\item H.~Lebesgue, \emph{Le\c cons sur l'int\'egration et la recherche des fonctions primitives}, Gauthier-Villars, 1904. Monotone functions and the fundamental theorem.
\item J.~Radon, ``Theorie und Anwendungen der absolut additiven Mengenfunktionen'', \emph{Sitzungsber. Akad. Wiss. Wien} 122 (1913), 1295--1438; O.~Nikodym, ``Sur une g\'en\'eralisation des int\'egrales de M.~J.~Radon'',
\emph{Fund. Math.} 15 (1930), 131--179.
\item G.~H.~Hardy and J.~E.~Littlewood, ``A maximal theorem with function-theoretic applications'', \emph{Acta Math.} 54 (1930), 81--116. The cricket story is on the first page.
\item G.~B.~Folland, \emph{Real Analysis}, 2nd ed., Wiley, 1999, Chapter~3. Our route through the maximal function and the differentiation of measures follows it.
\item E.~M.~Stein and R.~Shakarchi, \emph{Real Analysis}, Princeton, 2005, Chapter~3.
\item R.~Salem, ``On some singular monotonic functions which are strictly increasing'', \emph{Trans. Amer. Math. Soc.} 53 (1943), 427--439.
\end{itemize}
